\section*{Appendix C: Bullet Cluster Order-of-Magnitude Estimate}

We assess whether the DEV theory is qualitatively consistent with the
$\sim 200$~kpc lensing--gas centroid offset observed in 1E~0657-558.
At the core radius $r_c \approx 300$~kpc, the characteristic Newtonian
acceleration is $g_{\rm bullet} = G M_{\rm tot}/r_c^2 \approx 2.56e-09$~m\,s$^{-2}$,
giving $g_{\rm bullet}/a_0 \approx 21$. The Bullet Cluster therefore
operates deep in the Newtonian regime ($\mu \to 1$), where the DEV
modification of gravity is suppressed and the gravitational slip
$\eta - 1 \sim (2\beta/3)/\sqrt{x(1+x)}$ becomes small ($\sim 1.6e-03$ at $r_c$).
A simple ram-pressure lag model in which the gas is decelerated by
a fraction $f_{\rm decel} \in [0.1, 0.5]$ of $v_{\rm rel} \approx 3000$~km\,s$^{-1}$
over the merger crossing time predicts a centroid separation
$\Delta x \in [72, 360]$~kpc, bracketing the observed value.
We emphasize that this is an order-of-magnitude estimate; the actual
separation in DEV is governed by the inertia and rigidity of the vector
field $A_\mu$ during the merger, and a quantitative prediction requires
full numerical simulation, which we leave to future work.
